% ============================================================
%  CAPÍTULO 9: VALORACIÓN ÉTICA DEL PROYECTO
% ============================================================

\chapter{Capítulo Noveno: Valoración Ética del Proyecto}
\label{ch:etica}

Este capítulo recoge la valoración ética de \textit{TankGo}, identificando las cuestiones
éticas relevantes que plantea el proyecto, analizando las opciones disponibles y justificando
las decisiones adoptadas a lo largo del desarrollo. Se abordan cuatro dimensiones principales:
la privacidad y el tratamiento de los datos del usuario, la transparencia e interpretabilidad
del módulo de predicción basado en aprendizaje automático, el impacto social del sistema y
su alineación con los Objetivos de Desarrollo Sostenible (ODS) de las Naciones Unidas
\cite{onu_ods_goals}, y los principios de ética profesional del ingeniero aplicados
al proceso de desarrollo.

% ------------------------------------------------------------
\section{Privacidad y tratamiento de datos personales}
\label{sec:etica-privacidad}

\subsection{Identificación de la cuestión ética}

\textit{TankGo} maneja datos con implicaciones directas para la privacidad de sus usuarios.
La plataforma procesa: (i) la ubicación geográfica del usuario, necesaria para
mostrar estaciones próximas y calcular rutas de repostaje, (ii) las estaciones marcadas como
favoritas, que revelan patrones de movilidad y hábitos de consumo, y (iii) las consultas
realizadas al asistente de voz, que pueden contener información contextual sobre desplazamientos
y preferencias.

La cuestión ética principal es la tensión entre la
\textbf{utilidad de las funcionalidades personalizadas}, que requieren cierto grado de
persistencia o procesamiento de datos, y el \textbf{derecho a la privacidad} del usuario,
entendido como principio deontológico irrenunciable en el diseño de sistemas de información.

\subsection{Análisis de opciones}

Se han considerado dos aproximaciones principales para el tratamiento de estos datos:

\begin{itemize}
    \item \textbf{Opción A — Perfilado amplio:} almacenar el historial completo de consultas,
          ubicaciones y rutas para mejorar la personalización y las predicciones del modelo
          de \textit{machine learning}. Esta opción maximiza la utilidad funcional, pero compromete
          la privacidad, incrementa la superficie de ataque ante posibles vulnerabilidades y puede
          generar desconfianza en el usuario.

    \item \textbf{Opción B — Minimización de datos:} recopilar únicamente los datos
          estrictamente necesarios para cada funcionalidad concreta. La ubicación se procesa
          en tiempo real sin almacenamiento persistente más allá de la sesión, las predicciones
          de precio se calculan exclusivamente para las estaciones marcadas como favoritas por
          el usuario, de forma voluntaria y explícita.
\end{itemize}

\subsection{Decisión adoptada y justificación}

Se ha optado por la \textbf{Opción B}, basada en el principio de minimización de datos
(\textit{data minimisation}), coherente con el Reglamento General de Protección de
Datos (RGPD) \cite{rgpd2016}. Esta decisión no se deriva de un dilema de difícil
resolución, sino de la aplicación del principio de que \textbf{la privacidad prima por
defecto} cuando existe incertidumbre sobre la necesidad real de un dato. Desde el
principio de \textit{privacy by design} \cite{rgpd2016}, el sistema se ha diseñado para
no acumular datos que no sean imprescindibles para la funcionalidad ofrecida.

En la práctica, esto implica que las predicciones de LightGBM se generan y almacenan
únicamente para estaciones favoritas, reduciendo también el coste computacional, que
la geolocalización se solicita de forma explícita al usuario en cada sesión, y que el
asistente de voz no registra el historial de conversaciones más allá del contexto inmediato
de la consulta en curso.

% ------------------------------------------------------------
\section{Transparencia e interpretabilidad del modelo predictivo}
\label{sec:etica-ia}

\subsection{Identificación de la cuestión ética}

\textit{TankGo} incorpora un módulo de predicción de precios de carburante basado en
LightGBM con regresión por cuantiles. Este tipo de sistema plantea una cuestión ética
relacionada con la \textbf{confianza del usuario en las predicciones}: si el usuario
interpreta una predicción como un dato certero en lugar de como una estimación probabilística,
podría tomar decisiones de repostaje basadas en información que no tiene garantía de exactitud.

\subsection{Análisis de opciones}

\begin{itemize}
    \item \textbf{Opción A — Presentar las predicciones como valores exactos:} simplifica
          la interfaz, pero induce al usuario a sobrevalorar la fiabilidad del modelo.

    \item \textbf{Opción B — Comunicar explícitamente la naturaleza probabilística de las
              predicciones:} señalizar visualmente que se trata de estimaciones sujetas a incertidumbre,
          diferenciando los datos oficiales del MITECO \cite{datosgob_carburantes} de las
          predicciones generadas por el modelo.
\end{itemize}

\subsection{Decisión adoptada y justificación}

Se ha optado por la \textbf{Opción B}, con arreglo al principio de \textbf{autonomía
del usuario}: para que una decisión sea autónoma, el usuario debe contar con información
veraz sobre la naturaleza de los datos que se le presentan. La interfaz diferencia
visualmente los precios actuales, obtenidos en tiempo real de la fuente oficial, de las
predicciones del modelo, que se presentan con indicadores de incertidumbre derivados de
la estimación por cuantiles. Esta práctica está alineada con las recomendaciones de la
\textit{Trustworthy AI} propuesta por la Comisión Europea \cite{ec2019trustworthyai},
que incluye la transparencia como uno de los siete requisitos fundamentales que todo
sistema de IA debe satisfacer para considerarse digno de confianza.

% ------------------------------------------------------------
\section{Impacto social y alineación con los ODS}
\label{sec:etica-ods}

\textit{TankGo} nace con una vocación de utilidad social: facilitar al ciudadano
la comparación de precios de carburante y la localización de puntos de recarga eléctrica,
contribuyendo a una movilidad más informada y económicamente más justa.

Desde la perspectiva del impacto social, se identifica un riesgo menor aunque digno de
mención: la plataforma beneficia principalmente a usuarios con acceso a un dispositivo
móvil o navegador web y con vehículo propio, lo que podría suponer una brecha de
accesibilidad respecto a colectivos sin recursos tecnológicos o sin vehículo privado.
Sin embargo, dado que el servicio se ofrece de forma gratuita y accesible sin registro
obligatorio, se mitiga parcialmente esta asimetría.

% ------------------------------------------------------------
\section{Principios de ética profesional del ingeniero}
\label{sec:etica-profesional}

El desarrollo de \textit{TankGo} se ha guiado por tres principios fundamentales de la
ética profesional en ingeniería:

\textbf{Principio de beneficencia.} El sistema persigue un beneficio tangible para el
usuario final: reducir la brecha informativa del mercado de los carburantes,
permitiéndole tomar decisiones de repostaje más ventajosas. Asimismo, el fomento de
la movilidad eléctrica contribuye al bien común en términos medioambientales.

\textbf{Principio de autonomía.} La plataforma respeta y potencia la capacidad de
decisión del usuario. No se imponen recomendaciones, sino que se provee información
contrastada y herramientas de análisis (comparativa de precios, predicciones
etiquetadas como tales, recomendaciones de ruta explicadas) para que el usuario
tome sus propias decisiones de forma informada. Las funcionalidades personalizadas
(favoritos, predicciones específicas) son siempre de carácter voluntario y activadas
explícitamente por el usuario.

\textbf{Principio de justicia.} Los datos de precios utilizados proceden de fuentes
oficiales del Ministerio para la Transición Ecológica y el Reto Demográfico
\cite{datosgob_carburantes}, garantizando la objetividad e imparcialidad de la
información presentada. No se introducen mecanismos de promoción de estaciones
concretas ni se acepta ningún tipo de patrocinio que pudiera sesgar los resultados
mostrados al usuario.

% ------------------------------------------------------------
\section{Síntesis de la valoración ética}
\label{sec:etica-sintesis}

Las principales cuestiones éticas identificadas en \textit{TankGo} se articulan en torno
a la privacidad de los datos de movilidad del usuario, la transparencia en la comunicación
de las predicciones generadas por el modelo de aprendizaje automático, y el impacto social
de una herramienta orientada a democratizar el acceso a información sobre precios energéticos.
En todos los casos, las decisiones de diseño adoptadas priorizan la protección del usuario
frente a la maximización funcional, la transparencia frente a la opacidad algorítmica,
y el beneficio colectivo frente al interés particular. El proyecto contribuye
a una movilidad más sostenible, informada y económicamente
equitativa.